% ============================================================================
% ReSIDS v2.0 -- Monitoring and control interface
% Section for the v2 paper. \input{} it; requires booktabs only.
%
% Generated alongside docs/API.md, docs/COMMANDS.md and docs/MULTI_INSTANCE.md,
% which remain the normative documents; this is the paper-facing summary.
% ============================================================================

\section{Monitoring and control interface}
\label{sec:api}

A runtime-switchable IDS is only operable if the switching is observable and, where it
must be, authorised. This section specifies the interface through which an external
monitor observes ReSIDS and acts on it. We state it as a contract rather than as an
implementation, and we separate what is implemented from what is specified
(Section~\ref{sec:api-status}), because the distinction matters to anyone building against
it.

\subsection{Design principle: authority before protocol}
\label{sec:api-authority}

The interface follows from a single architectural decision: \emph{autonomy is granted only
in the direction that fails safe}. Table~\ref{tab:authority} lists the four situations in
which the system changes shape. Exactly one of them is autonomous.

\begin{table}[htbp]
\centering
\caption{Authority rule. One autonomous action; everything else is monitor-governed. The
asymmetry is driven by measurement: losing nodes in GL costs nothing, losing them in FL
costs F1.}
\label{tab:authority}
\begin{tabular}{@{}llll@{}}
\toprule
Situation & Authority & Latency & Fallback \\
\midrule
Node inactive while in FL $\rightarrow$ GL & \textbf{ReSIDS, autonomous} & $\sim$2 rounds & deadline $D$ \\
Further node loss while in GL & monitor & policy & none \\
Node with attributed intrusion & monitor & policy & none \\
GL $\rightarrow$ FL return & \textbf{monitor} & policy & \textbf{none, by design} \\
\bottomrule
\end{tabular}
\end{table}

In GL, recall is unaffected by node loss---it holds at 100\% from 14 nodes down to 3,
because the retained booster union keeps every specialist---so there is no urgency and no
reason to spend autonomy. In FL the same loss removes a specialist from the aggregate and
degrades detection (96.02 $\rightarrow$ 86.33 F1 under cascading loss), and the aggregator
is a single point of failure. Being stuck in FL with a dead aggregator is therefore
dangerous and the transition out of it must not depend on an external party; being stuck in
GL is merely slower ($\Theta(N)$ dissemination instead of depth-1), so the return to FL
deliberately has \emph{no} fallback. If the monitor is unreachable, the system rests in the
degraded-but-safe state.

Returning to FL means re-accepting the hub---the very point of failure and of trust that GL
exists to survive---which is a trust decision rather than a heuristic one. Consequently the
recovery gates (dwell plus unanimity of local health) cease to be an autonomous commit rule
and become \emph{evidence} the agent reports and checks a command against.

Every event carries \texttt{decided\_by} $\in$ \{\texttt{autonomous}, \texttt{monitor},
\texttt{operator}\}, present only on events that \emph{are} an action. This makes the rule
machine-checkable rather than documentary: a conformance script rejects any stream that
claims autonomy where Table~\ref{tab:authority} does not grant it.

\subsection{Observation surface}
\label{sec:api-events}

Five event types (Table~\ref{tab:events}) carry every state change an operator needs. Two
are actions; three are observations. An event refers to a closed traffic window and carries
its bounds, so it can be correlated with sequence-of-event records: \texttt{window\_start}
and \texttt{window\_end} are the window, distinct from \texttt{ts}, the emission instant.

\begin{table}[htbp]
\centering
\caption{Event types. Only the two actions carry \texttt{decided\_by}; attaching it to an
observation would assert a decision that was never taken.}
\label{tab:events}
\begin{tabular}{@{}llp{0.42\linewidth}@{}}
\toprule
Type & Kind & Payload beyond the common fields \\
\midrule
\texttt{architecture\_change} & action & \texttt{from\_mode}, \texttt{to\_mode}, \texttt{reason}, \texttt{decided\_by} \\
\texttt{node\_isolated} & action & \texttt{source\_node}, \texttt{reason}, \texttt{decided\_by} \\
\texttt{node\_failure} & observation & \texttt{failed\_nodes} \\
\texttt{node\_recovery} & observation & \texttt{recovered\_nodes} \\
\texttt{intrusion\_detected} & observation & \texttt{attack}, \texttt{k\_votes}, \texttt{n\_flags}, \texttt{detector\_specialists}, \texttt{source\_node} \\
\bottomrule
\end{tabular}
\end{table}

Two naming decisions encode architectural facts. \texttt{detector\_specialists} names
boosters, not nodes: every node holds the same diffused union and reaches the same verdict,
so a per-node reading would suggest a corroboration the architecture does not have.
\texttt{n\_flags} is the alarm's \emph{volume} (flagged samples in the window) while
\texttt{k\_votes} is its \emph{corroboration strength} (concurring specialists on the
strongest single sample); Section~\ref{sec:api-fp} shows why both must be reported.

\subsubsection{Behavioural contract}

What a consumer may rely on is not expressible in a schema, so it is stated separately:
events are ordered by \texttt{seq} ascending; \texttt{seq} starts at 1, has no gaps, and
continues across restarts; \texttt{since} is exclusive; retention is bounded and a
\texttt{since} below it yields \texttt{410 Gone} rather than a silently short answer.

The operative rule is that a consumer must reconcile using the highest \emph{contiguous}
\texttt{seq} it has received, never the highest \emph{seen}. The difference is not
cosmetic. Once a later notification arrives, the tail is current while an earlier event is
still missing, and the hole becomes invisible. Measured on the replay harness at 20\%
notification loss: \textbf{33 events permanently lost} with the tail test, \textbf{zero}
with the contiguous one. The same request subsumes liveness---an empty answer means alive,
no gap, and nothing to recover---so a separate periodic probe costs a round trip and buys
nothing.

\subsubsection{False positives are the monitor's design constraint}
\label{sec:api-fp}

The interface must be sized for the alarm load it will actually carry. Replaying a
benign-only stream, in which every alarm is by construction a false positive, 252 of 724
windows raise one (34.8\%). The two populations separate on volume rather than on
corroboration: median \texttt{n\_flags} of 56 for false positives against 973 for true
alarms. A threshold of $\texttt{n\_flags} \geq 275$ removes every false positive while
retaining 91.8\% of true alarms; the remaining 8.2\% fall below the worst false positive on
volume and are recoverable only through \texttt{k\_votes}, whose range overlaps
($\leq 4$ on false positives, up to 10 on true ones). Reporting a single fused score
instead of both fields would make that residue unrecoverable, which is the reason the event
carries two numbers and no verdict.

\subsection{Command surface}
\label{sec:api-commands}

Three commands---and no more---cover every row of Table~\ref{tab:authority} that is not
autonomous: \texttt{set\_mode} for both transitions, \texttt{isolate\_node} and
\texttt{readmit\_node} for membership. Readmission is not symmetry for its own sake: with a
34.8\% false-positive rate on benign windows, wrongful isolation is expected, and a surface
that can isolate but not readmit would force a node restart to undo the IDS's own mistake.

Four envelope fields carry the semantics, each answering a failure mode rather than a
convention. \texttt{command\_id}, a client-generated UUID, makes execution \emph{at most
once}: the channel that loses messages also duplicates them, and a duplicated isolation or
mode switch is not idempotent by nature. \texttt{authority\_epoch}, monotonic, refuses a
commander whose epoch has been superseded, without which two monitors during a failover are
indistinguishable from legitimate traffic. \texttt{not\_after} bounds replay of a captured
command. \texttt{expect} states the view the monitor decided on, and a mismatch refuses the
command outright---the same argument as contiguous-seq reconciliation, one layer up: a
monitor decides on a polled view that is already in the past, and an actuation surface must
not let a stale view act.

Three properties are not negotiable.

\begin{enumerate}
\item \textbf{The deadline $D$ is not commandable.} No verb cancels, holds or extends it. A
commander able to postpone $D$ is a commander able to keep a node in FL with a dead
aggregator, reintroducing through the control channel the single point of failure GL exists
to survive. $D$ is configuration, changed by redeploying policy, not by a message in
flight.

\item \textbf{The monitor decides \emph{whether}; the agent verifies whether it is
\emph{possible}.} A GL $\rightarrow$ FL return is refused while membership is incomplete,
and an isolation is refused below $\texttt{min\_nodes} = 3$---the floor at which the
measurement stops supporting the claim, since recall is verified from 14 nodes down to 3 and
no further. An explicit \texttt{force} flag overrides both and is recorded in the resulting
event.

\item \textbf{The response is not the commit record.} A response can be lost; the
authoritative record is the event, carrying \texttt{decided\_by} and \texttt{command\_id}.
Rejected commands produce no event, so every answered command also receives a position in a
command log that reconciles exactly like the event stream---otherwise a refused actuation
attempt would leave no trace at all.
\end{enumerate}

Isolation removes a node from the federation and the booster pool; it does \emph{not}
disconnect it from the process bus. An isolated node keeps publishing GOOSE/SV and keeps
being attributed as the source of alarms. Containment is a network action that ReSIDS
neither performs nor claims, and the command result says so on every accepted isolation.

What the surface deliberately omits is as much a part of the specification as what it
offers: no booster upload (a commander able to inject a model can poison detection), no
runtime change of $k$ or of the decision thresholds (a silent way to blind the IDS), no
supply of labels or training data (a commander able to label traffic can teach the system
that an attack is benign), and no traffic blocking.

\subsection{Deployment: fifteen instances}
\label{sec:api-instances}

In the decentralised deployment each agent serves the interface independently: the FL
aggregator plus the 14 clients, fifteen endpoints in all. Four things change relative to a
single endpoint, and three of them are places where a uniform treatment of the fifteen is
wrong.

\paragraph{Fifteen watermarks.} Each instance has its own \texttt{seq} space, so a
\texttt{since} from one endpoint is meaningless on another. The monitor holds one contiguous
watermark per instance; a shared counter would be silently wrong rather than merely lossy.

\paragraph{The aggregator's silence is not a fault.} In GL the aggregator has no role and
stops emitting. A monitor treating its fifteen endpoints uniformly declares it dead
immediately after the FL $\rightarrow$ GL switch it has just observed---that is, precisely
when the operator most needs an unambiguous picture. The rule is derived from the stream
rather than special-cased: while the mode is \texttt{gossip}, the aggregator is expected
quiet.

\paragraph{Agreement is not corroboration.} Alarms for one window arrive from many
instances, and what that means depends on what the instances observe. When every agent
scores the whole window with the same diffused union, all live agents reach the same
verdict \emph{by construction}; fourteen identical alarms are agreement, and counting them
as corroboration would inflate confidence with replication. The useful signal there is the
converse---divergence---and an isolated-but-live node is exactly that case, its membership
view frozen at the moment of isolation while its agent continues. When instead each agent
observes its own bus segment, alarms from different instances are independent evidence and
$k$-of-$n$ fusion \emph{at the monitor} becomes a second, outer layer above the $k$-of-$n$
already running inside each node's pool. Section~\ref{sec:api-corrob} quantifies the
difference. We report both because only the first is faithful to the single-trace
evaluation used throughout this paper; the second is a declared modelling premise.

\paragraph{Fan-out, with asymmetric unanimity.} A mode switch is federation-wide while
commands are per-instance, so a fan-out can partially fail---and the two directions do not
tolerate it equally. A partial FL $\rightarrow$ GL is safe: the stragglers converge on GL,
the resting state. A partial GL $\rightarrow$ FL splits the federation, leaving some
instances aggregating through a hub the others have abandoned. The return therefore requires
unanimity and is aborted otherwise, while the outbound switch does not. The fifteen commands
of one decision share an \texttt{intent\_id} so the audit trail records them as a single
act, each retaining its own \texttt{command\_id}. This is the only place where the
multi-instance setting changes command semantics rather than repeating them, and it follows
from the same asymmetry as Table~\ref{tab:authority}.

\subsection{Measured: agreement is not evidence}
\label{sec:api-corrob}

We instantiated the deployment of Section~\ref{sec:api-instances}---fifteen endpoints served
by the reference implementation, polled over HTTP by the monitor---and measured how many
instances report the same traffic window, on attack traffic and on benign-only traffic in
which every alarm is by construction a false positive. The monitor consumed the streams with
no gap on any instance, each reconciled on its own contiguous sequence number.
Table~\ref{tab:corrob} gives the result.

\begin{table}[htbp]
\centering
\caption{Client instances reporting the same window. Requiring corroboration across
instances costs 6.4\% of true detections and filters no false positive at all: a false
positive is a property of the traffic window, not an idiosyncrasy of an agent.}
\label{tab:corrob}
\begin{tabular}{@{}llrr@{}}
\toprule
Window class & View & $\geq 2$ inst. & $\geq 5$ inst. \\
\midrule
Attack ($n=282$) & \texttt{replicated} & 100.0\% & 100.0\% \\
                 & \texttt{sharded} & \textbf{93.6\%} & 93.3\% \\
\addlinespace
Benign ($n=252$), all false positives & \texttt{replicated} & 100.0\% & 100.0\% \\
                 & \texttt{sharded} & \textbf{100.0\%} & 100.0\% \\
\bottomrule
\end{tabular}
\end{table}

The replicated rows are uninformative by construction, and that is itself the first result.
Fourteen identical alarms are not fourteen corroborations: every agent holds the same
diffused union and scores the same samples, so a monitor fusing them as independent evidence
would report a confidence never independently obtained.

The two emphasised figures are the second and more consequential result:
\textbf{instance-level corroboration is strictly harmful as a filter}. Requiring $k \geq 2$
\emph{instances} discards 6.4\% of true detections while removing no false positive
whatsoever---every false positive is corroborated by at least seven instances even when each
agent observes only its own slice. The reason is structural: a false positive is a property
of the traffic window rather than an idiosyncrasy of an agent, so a benign window that
resembles an attack resembles one to everybody, and partitioning the observation does not
decorrelate the error. It only reduces how much of it each agent sees. Triage must
therefore remain volume-based, as the single-endpoint baseline of Section~\ref{sec:api-fp}
indicated. Fanning out to fifteen endpoints buys reconciliation, attribution and divergence
detection; it does not buy a second opinion.

On the attack side the sharded distribution is bimodal---either all fourteen instances fire
or exactly one does---and the single-instance tail consists of the sparse windows, in which
so few samples fall inside the second that only one agent's slice contains any. One artifact
must be declared: those windows concentrate on a single instance, because round-robin slice
assignment gives index 0 the first sample of every window. A deployment partitioned by bus
segment would not carry that bias, so the 6.4\% figure stands while its attribution to one
agent does not.

\paragraph{Divergence as a detector.} Because agreement is uninformative under replication,
the useful signal is its absence. Of the five nodes isolated during the scenario, the
monitor identified \textbf{four from divergence alone}, without reading a single
\texttt{node\_isolated} event: each isolated instance continues to report, carrying the
membership view frozen at the moment it was cut out of the federation, which places it in
the minority against the others. No instance ever disagreed about \emph{which} attack was
present; the divergence was entirely in membership. The fifth isolation is invisible to this
signal for a structural reason worth stating---it froze at what was already the final
membership, so no subsequent change exposed the staleness. Divergence therefore detects
isolation only while membership continues to change, and cannot be the sole mechanism.

\subsection{Security}
\label{sec:api-security}

The observation surface is read-only and, bound to loopback, unauthenticated. That stance
does not extend to the command surface, which changes the behaviour of a protection system:
mutual authentication is mandatory, and each command additionally carries a detached
signature. The reason for the second mechanism is that channel authentication proves the
\emph{connection}, not the \emph{decision}: with a TLS-terminating proxy or an IoT hub in
the path---which a CoAP deployment explicitly introduces---that intermediary becomes an
implicit commander, and an audit trail based on channel identity alone cannot tell.

Regarding standards alignment, IEC 62351 offers no DTLS profile; it delegates TCP/IP
profiles to TLS (62351-3) and protects GOOSE/SV inside the PDU (62351-6, 62351-9). The
ReSIDS-to-monitor channel is out-of-band management traffic and therefore outside that
scope, which we state explicitly rather than claim conformance by proximity. Transport cost
is not the deciding factor at our volume: DTLS 1.3 adds roughly 11 bytes per record, which
still fits inside Ethernet minimum padding, so a protected datagram occupies the same 84
bytes on the wire as an unprotected one.

The command surface creates an attack surface that did not exist while the monitor was
read-only. A compromised commander can force GL, force FL, isolate honest nodes, or readmit
an isolated one. The mechanisms above are bounds rather than solutions---epoch for failover,
membership floor, flap limit, signature for attribution, and the deliberate absence of a
deadline veto so that the worst case remains \emph{degraded} rather than \emph{disarmed}.
Byzantine-aware trust of the commander itself is outside the scope of this version.

\subsection{Implementation status}
\label{sec:api-status}

\begin{table}[htbp]
\centering
\caption{What is implemented, and what is specified. The separation is stated here because
an interface paper that blurs it is not actionable.}
\label{tab:api-status}
\begin{tabular}{@{}lp{0.52\linewidth}@{}}
\toprule
State & Component \\
\midrule
Implemented & event emission; \texttt{/events}, \texttt{/status}, \texttt{/health}; the
JSONL audit trail; contiguous-seq reconciliation in the reference consumer; the fifteen
endpoints and the multi-instance monitor \\
Specified only & \texttt{/commands} and the command log; at-most-once execution; mutual
authentication, per-command signature and epoch enforcement; \texttt{switch\_pending} and
the evidence fields of \texttt{/status}; the monitor-commanded switch policy and deadline
$D$ \\
\bottomrule
\end{tabular}
\end{table}

The event half of the loop is verifiable today: both scenario streams conform to the schema
and to the authority rule, and the command log those streams \emph{imply} is derived
mechanically from their monitor-decided events---seven commands for seven actions---which
demonstrates that the command surface is sufficient for the scenarios without claiming it is
built.
